
\chapter*{第五版前言}


华东师范大学数学系编写的《数学分析》(上、下册) 自 1980 年初版诞生以来,相继于 1990 年和 2001 年两次再版, 1987 年荣获全国第一届高等学校优秀教材优秀奖, 2004 年荣获上海市优秀教材一等奖,并于 2010 年出版第四版。自第四版出版以来,教材向着信息化和精美化的方向发展, 在使用中也发现了一些还可以继续完善的地方, 在此背景下我们着手第五版修订。

为做好本次修订,前期我们进行了广泛调研,征集使用意见,并召开了教材修订会,而后根据当前教材发展的新形势,最终确定第五版的修订方案。本次修改的主要原则是基本保留第四版的风格,对某些不完善之处进行必要的处理,并适当补充数字资源(以图标 \(\times\) , 手机 示意)。

本次修改的主要内容有:

1. 增加了对函数 \(y = {x}^{n},x > 0\) 值域的讨论,从而明确了其反函数 \(y = \sqrt[n]{x}\) 的定义域,这就补上了第四版的一个不足;

2. 为了保证大学数学与中学数学内容的衔接, 根据授课老师的建议, 此次再版将三角函数的和差化积公式、积化和差公式作为习题出现在教材上；

3. 在有理函数的部分分式分解讲解中,我们适当作了一些说明；

4. 对 \(y = {\left( 1 + x\right) }^{\alpha }\) 幂级数的展开式处理稍作修改;

5. 对一些例题进行适当的增减;

6. 将“微积分学简史”作为数字资源。

参加第五版编写工作的老师有:

庞学诚任主编,并负责编写第一章至第七章；

吴 畏 负责编写第八章至第十一章,以及第十三章；

柴 俊 负责编写第十二章、第十四章至第十八章；

戴浩晖负责编写第十九章至第二十三章。

在编写过程中,一直得到我校数学科学学院老师的关心和帮助,高等教育出版社的编辑也付出了辛勤的劳动,对此我们表示衷心的感谢。在这里我们还要感谢复旦大学楼红卫教授、浙江大学阮火军教授,同时还要感谢在本书的修订和出版过程中积极提出建议的所有朋友。

本书自 1980 年第一版出版以来,深受广大读者的关爱与支持,在此我们一并致以深切谢意！并衷心希望读者在阅读和使用本教材的过程中继续提出宝贵意见。

编 者

2018 年 2 月

\newpage

\chapter*{提示}

本排版是对用户hnkznhb@126.com的初排本进行的二次排版。修饰了一些细节问题,本讲义仅供排版练习，所有版权属于原作者.
